
This work addresses a gap at the intersection of MCP-based tool composition, automated validation, and semantic reasoning analysis. Four categories of prior work provide context.

\textbf{MCP Tool Composition Frameworks.} Ahn et al. (2025) introduced an agentic Model Context Protocol framework for medical concept standardization, demonstrating zero-training validation without hallucination by mapping clinical terms to standardized ontologies using MCP to integrate external knowledge bases with LLM reasoning. This work extends Ahn et al. in three ways. First, it applies MCP to research pipelines rather than medical domains---research workflows have messier phase boundaries and evolving hypotheses compared to standardized medical concepts. Second, it treats MCP traces as validation artifacts rather than tool composition logs. Third, it empirically validates that 97.48\% of research pipeline tool calls contain natural language in both queries and results, establishing MCP traces as rich semantic artifacts beyond structured medical data. Phase 1 literature analysis found that only 1 of 15 recent ML infrastructure papers uses MCP explicitly, revealing a research gap in MCP-based semantic validation.

\textbf{Agent-Driven Annotation and Automation.} Fu et al. (2025) introduced PRDBench, an agent-driven pipeline for automated benchmark construction that reduces annotation costs by over 90\% while maintaining >90\% human alignment using LLM agents to generate diverse project-level programming tasks from GitHub repositories. This work complements Fu et al. by achieving zero-annotation validation via trace analysis instead of reducing annotation cost via generation. Where Fu et al. generate benchmarks requiring human verification, the framework presented here infers validation criteria from execution traces without any manual labeling, achieving 82.7\% recall and 86.3\% precision using only zero-shot LLM prompts and multi-vote consensus. The key difference: Fu et al. reduce annotation overhead for benchmark creation, this work eliminates annotation for pipeline validation by exploiting MCP trace structure.

\textbf{Declarative Constraint Enforcement.} Neutatz et al. (2021) demonstrated that declarative feature selection can satisfy multiple constraints simultaneously (fairness, privacy, execution time) in ML systems through explicit constraint declaration and multi-objective optimization. This work diverges fundamentally in constraint specification method. Neutatz et al. require manual declaration of constraints, while the approach here infers constraints from execution traces. This distinction matters for research pipelines where constraints evolve across phases and may not be explicitly stated upfront. However, the negative result on constraint inference reveals a limitation of inference-based approaches: inferred constraints face semantic ambiguity that declarative constraints avoid through precise verification.

\textbf{Traditional Pipeline Validation Tools.} Production ML pipeline tools like MLflow, DVC, and Great Expectations provide experiment tracking, versioning, and data quality validation but require manual test writing and lack semantic reasoning capabilities. These tools are MCP-agnostic, operating on files, metrics, and schemas without understanding tool-calling semantics. The framework presented here is MCP-native, exploiting the structure of tool calls to extract semantic content. The trade-off: traditional tools offer battle-tested reliability for syntactic validation with high precision and low false positive rates, while semantic approaches introduce LLM-based uncertainty. These are complementary: syntactic validation leverages established tools, semantic layers address the reasoning gap traditional tools miss.

\textbf{Positioning Summary.} The unique contribution is treating MCP traces as semantic artifacts for validation-as-inference rather than validation-as-specification. This extends Ahn et al. by validating traces (not orchestrating tools), complements Fu et al. by achieving zero-annotation (not reducing annotation), diverges from Neutatz et al. by inferring constraints (not declaring them), and augments traditional tools by adding semantic layers atop syntactic validation.
